\documentclass[11pt]{article}
\usepackage{amsmath,amssymb,amsthm,mathtools,enumitem,booktabs}
\usepackage[margin=1in]{geometry}

\newtheorem{theorem}{Theorem}
\newtheorem{lemma}{Lemma}
\newtheorem{proposition}{Proposition}
\newtheorem{definition}{Definition}
\newtheorem{warning}{Warning}
\newtheorem{corollary}{Corollary}

\newcommand{\Kcmp}{\mathsf K_{\rm cmp}}
\newcommand{\AZ}{\mathsf A_Z}
\newcommand{\EE}{E_{\rm EF}}
\newcommand{\Ccore}{\mathcal C^-_{\log}}
\newcommand{\ip}[2]{\langle #1,#2\rangle}

\title{Step 76: Signed-to-Positive Matching for the Completed Weil Feature Gate}
\author{Six Birds / Formed-Layer Membrane Program}
\date{}

\begin{document}
\maketitle

\section*{Purpose}
Step 75 proposed a log-side anti-invariant core
\[
  \Ccore=C_c^\infty(\mathbb R)^{J=-1},\qquad (Jf)(t)=\overline{f(-t)},
\]
and a completed candidate feature map
\[
  W_{\rm cmp}=\begin{bmatrix}W_{\rm pr}\\ W_\infty\\ W_{\rm pole}\\ W_{\rm tail}\end{bmatrix}.
\]
Step 76 asks what it would mean for these features to match the completed Weil explicit formula strongly enough to attempt the Douglas bridge
\[
  V_Z=T W_{\rm cmp},\qquad \|T\|\le 1,
\]
or equivalently
\[
  \AZ=V_Z^*V_Z\preceq W_{\rm cmp}^*W_{\rm cmp}=\Kcmp.
\]
The main point is that a signed explicit formula is not yet a positive carrier.  It becomes one only after a positive feature representation, a form-core/closability record, and a domination record.

\section{Signed Weil data and positive feature data}

\begin{definition}[Signed Weil datum on a core]
A signed Weil datum on a complex vector space/core \(\mathcal C\) consists of Hermitian quadratic forms
\[
 q_{\rm pr}^{\rm raw},\quad q_\infty^{\rm raw},\quad q_{\rm pole}^{\rm raw},\quad q_{\rm tail}^{\rm raw},\quad q_Z
\]
with the intended identity
\[
 q_W[f]=q_{\rm pr}^{\rm raw}[f]+q_\infty^{\rm raw}[f]+q_{\rm pole}^{\rm raw}[f]+q_{\rm tail}^{\rm raw}[f]-q_Z[f].
\]
The form \(q_Z\) is the zero-side anti-invariant displacement form.
\end{definition}

\begin{definition}[Positive completed feature realization]
A positive completed feature realization of the signed datum is a Hilbert feature space
\[
 \mathcal E_{\rm cmp}=\mathcal E_{\rm pr}\oplus\mathcal E_\infty\oplus\mathcal E_{\rm pole}\oplus\mathcal E_{\rm tail}
\]
and linear maps
\[
 W_i:\mathcal C\to \mathcal E_i,
 \qquad V_Z:\mathcal C\to\mathcal Z,
\]
such that
\[
 q_{\rm cmp}[f]=\sum_i\|W_i f\|^2,
 \qquad q_Z[f]=\|V_Zf\|^2,
\]
and the signed Weil form is represented as
\[
 q_W[f]=q_{\rm cmp}[f]-q_Z[f]
 \qquad (f\in\mathcal C).
\]
\end{definition}

\begin{theorem}[Weil positivity to Douglas gate]
Suppose \(q_{\rm cmp}\) and \(q_Z\) have a positive completed feature realization on a dense form core \(\mathcal C\), and suppose \(q_{\rm cmp}\) is closable with \(\mathcal C\) as a form core for the completed response space.  If
\[
 q_Z[f]\le q_{\rm cmp}[f]
 \qquad (f\in\mathcal C),
\]
then the inequality extends to the completed response space and there is a contraction
\[
 V_Z=T W_{\rm cmp},\qquad \|T\|\le 1.
\]
Conversely, such a contraction implies \(q_Z\le q_{\rm cmp}\).
\end{theorem}

\begin{proof}
On the core, \(\|V_Zf\|^2\le\|W_{\rm cmp}f\|^2\).  Hence the rule
\[
 W_{\rm cmp}f\mapsto V_Zf
\]
is well-defined on \(\operatorname{Ran}(W_{\rm cmp}|_{\mathcal C})\), because if \(W_{\rm cmp}f=0\) then \(V_Zf=0\). It is contractive on this range and therefore extends by continuity to a contraction on the closure. Extending it by zero on the orthogonal complement gives \(T\). The converse follows by applying \(\|T\|\le1\).
\end{proof}

\section{What must be matched}

\begin{proposition}[Signed-to-positive matching gate]
A classical signed explicit formula can enter the membrane framework only through one of the following accepted records:
\begin{enumerate}[label=(\alph*)]
\item an exact positive feature realization \(q_W=q_{\rm cmp}-q_Z\) and \(q_W\ge0\) on the full completed core;
\item a defected realization
\[
 q_Z[f]\le q_{\rm cmp}[f]+e_{\rm EF}[f],
\]
with a positive explicit-formula defect \(e_{\rm EF}\);
\item a scoped nonclaim recording that the signed identity is only support evidence and does not certify the Douglas bridge.
\end{enumerate}
Trace equality, scalar prime-zero balance, or positivity on a finite test family alone does not pass this gate.
\end{proposition}

\begin{warning}[No free positiveization]
An indefinite Hermitian form cannot be represented as \(\|Wf\|^2\). It can only be represented as a difference of positive forms, repaired by adding positive completion terms, or recorded as a defect. Therefore the signed prime/gamma/pole/tail presentation of the explicit formula is not itself a completed carrier.
\end{warning}

\section{Translation-invariant shift features}

For a positive even measure \(\mu\) on \(\mathbb R\), define
\[
 (W_\mu f)(a,t)=f(t+a)-f(t).
\]
Then
\[
 q_\mu[f]=\int\!\int |f(t+a)-f(t)|^2\,dt\,d\mu(a)
 =\int \Psi_\mu(\xi)|\widehat f(\xi)|^2\,d\xi,
\]
where
\[
 \Psi_\mu(\xi)=\int |e^{ia\xi}-1|^2\,d\mu(a)
 =2\int (1-\cos(a\xi))\,d\mu(a)\ge0.
\]

\begin{definition}[Shift-feature cone]
A nonnegative even symbol \(\Psi\) belongs to the shift-feature cone if
\[
 \Psi(\xi)=2\int (1-\cos(a\xi))\,d\mu(a)
\]
for some positive even measure \(\mu\) satisfying the integrability needed for \(q_\mu\) on \(\Ccore\).
\end{definition}

\begin{proposition}[Translation-invariant positive feature criterion]
A translation-invariant quadratic form
\[
 q[f]=\int \Psi(\xi)|\widehat f(\xi)|^2\,d\xi
\]
is a positive shift-difference feature form exactly when \(\Psi\) lies in the shift-feature cone. In particular, positivity of \(\Psi\) is necessary but not by itself an accepted shift-feature record; the measure representation and form-domain record are part of the gate.
\end{proposition}

\section{Candidate slot statuses}

\subsection*{Prime-power slot}
The candidate
\[
 (W_{{\rm pr},L}f)_n(t)
 =\sqrt{\frac{\Lambda(n)}{\sqrt n}}\,\alpha_L(\log n)(f(t+\log n)-f(t))
\]
is a positive shift feature at every finite cutoff. It contributes symbol
\[
 \Psi_{{\rm pr},L}(\xi)=2\sum_{n\ge2}\frac{\Lambda(n)}{\sqrt n}\alpha_L(\log n)^2(1-\cos(\xi\log n)).
\]
Unearned record: exact matching to the prime side of the completed explicit formula, including cutoff/tail law.

\subsection*{Archimedean slot}
The candidate kernel
\[
 \kappa_\infty(a)\sim \frac1{2\sinh(|a|/2)}
\]
formally defines a positive shift-difference feature because the difference factor cancels the near-zero singularity:
\[
 |f(t+a)-f(t)|^2\sim a^2|f'(t)|^2,
 \qquad a^2|a|^{-1}\text{ is locally integrable.}
\]
Unearned record: exact renormalized matching to the gamma factor, including constants, null modes, and pole interactions.

\subsection*{Pole/completion slot}
Finite rank features
\[
 W_{\rm pole}f=(\langle f,p_1\rangle,\ldots,\langle f,p_m\rangle)
\]
are positive. Unearned record: exact declaration of pole/null directions and proof that they repair the signed completion rather than hide it.

\subsection*{Tail/support slot}
This slot records omitted prime shifts, support leakage, and finite-to-completed promotion.  It is accepted either as a positive feature or as an explicit defect \(E_{\rm tail}\) participating in a fixed/exhaustive ledger squeeze.

\section{Countermodels}

\begin{proposition}[Trace balance is not domination]
There exist positive matrices \(A,K\succeq0\) with \(\operatorname{tr}A=\operatorname{tr}K\), but \(A\npreceq K\).
\end{proposition}

\begin{proof}
Take
\[
 A=\begin{pmatrix}2&0\\0&0\end{pmatrix},\qquad
 K=\begin{pmatrix}1&0\\0&1\end{pmatrix}.
\]
Both have trace \(2\), but \(K-A=\operatorname{diag}(-1,1)\) is indefinite.
\end{proof}

\begin{proposition}[Finite core support is not full bridge]
Positivity of \(q_{\rm cmp}-q_Z\) on a proper finite-dimensional subspace does not imply positivity on the completed response space.
\end{proposition}

\begin{proof}
Let \(Y=\mathbb R^2\), \(\mathcal C=\operatorname{span}(e_1)\), and
\[
 K-A=\begin{pmatrix}1&0\\0&-1\end{pmatrix}.
\]
The form is positive on \(\mathcal C\) but not on \(Y\).
\end{proof}

\section{Conclusion}
Step 76 turns the matching problem into a precise gate:
\[
 \text{signed explicit formula}
 \longrightarrow
 \text{positive completed feature representation}
 \longrightarrow
 q_Z\le q_{\rm cmp}
 \longrightarrow
 \text{Douglas contraction}.
\]
The prime shift slot is already a positive feature candidate. The archimedean slot is plausible as a positive renormalized shift feature. The pole/completion and tail slots are necessary to avoid fake null budgets and moving-ledger overreads. The missing work is exact matching to the completed Weil formula on a completed anti-invariant response space.

\end{document}
